\section{Introduction}\label{sec:introduction}

Benchmark scores are the primary currency for comparing language models, yet training data contamination undermines the validity of these comparisons \citep{hendrycks2021mmlu,chang2025benchmarking}. When models encounter test items during training, observed scores conflate genuine capability with memorization, and practitioners have no principled way to determine how much of a reported accuracy reflects real understanding. The problem is compounded by the scale of modern training corpora: trillion-token web crawls make contamination pervasive and instance-level measurement infeasible \citep{oren2024proving}. Post-hoc calibration tools that quantify contamination effects despite imperfect detection are therefore essential for trustworthy evaluation.

Two developments make this gap urgent. First, \citet{sun2025emperor} systematically evaluate twenty contamination mitigation strategies (rephrasing, option shuffling, dynamic benchmark regeneration, and others) and demonstrate that none restores clean performance after contamination has occurred. This result forecloses the ``prevent contamination'' approach for models already trained on unrestricted web data. Second, contamination-free benchmarks \citep{white2024livebench} offer a partial solution but sacrifice the comparability and institutional knowledge accumulated around established benchmarks. Dynamic regeneration of test items may also be contaminated within weeks as web crawlers index benchmark repositories \citep{dong_2024_cdd}. The field therefore needs post-hoc statistical tools that can separate genuine ability from contamination artifacts on existing benchmarks.

Detection methods such as Min-K\%++ \citep{zhang_2024_minkpp}, ConStat \citep{dekoninck2024constat}, and Kernel Divergence Signatures \citep{choi_2025_kds} can flag contaminated items but stop at detection. They do not quantify how much scores are inflated or produce adjusted scores with uncertainty estimates. Score adjustment methods like DCR \citep{xu2025dcr} provide point corrections without uncertainty quantification or a principled measurement framework. IRT-based evaluation frameworks \citep{polo_2024_tinybenchmarks,habba_2025_growing_pains} apply psychometric ability estimation to improve evaluation efficiency, but they treat all items as equally trustworthy and ignore contamination entirely \citep{lalor_2024_irt_nlp}. These two research threads have developed in isolation: detection papers end at binary flags, while IRT papers assume clean data. No existing method bridges detection signals and a measurement-theoretic calibration framework to produce decontaminated ability estimates with credible intervals.

We propose ContamCal, a two-stage framework that connects contamination detection to IRT-based calibration. Stage~1 fuses multiple detection signals into continuous per-item exposure scores via training-free Reciprocal Rank Fusion \citep{cormack_2009_rrf}, avoiding the meta-circularity of learned aggregators that would require the very contamination labels we lack. Stage~2 embeds these exposure scores as covariates in a contamination-augmented IRT model (2PL+$\gamma$) that jointly estimates latent ability $\theta$ and a contamination effect parameter $\gamma$. The key empirical finding is that IRT ability estimation itself provides the dominant calibration benefit: vanilla 2PL without any contamination parameter reduces calibration error on ARC by 53\% (MAE from 0.034 to 0.016), comparable to the full 2PL+$\gamma$ model. The $\gamma$ parameter serves a complementary diagnostic role, providing directional information about where contamination effects are positive or negative across cognitive domains, even when estimates are underpowered for individual domain significance at small ensemble sizes.

Our contributions are as follows:
\begin{enumerate}
    \item \textbf{IRT calibration decomposition.} We demonstrate that IRT ability estimation provides calibration benefit independent of whether contamination is explicitly modeled through $\gamma$. This value decomposition clarifies that the primary benefit of psychometric modeling for contaminated benchmarks lies in separating stable ability from item-level noise, not in the contamination correction itself.
    \item \textbf{Multi-signal exposure detection via RRF.} We develop a training-free signal fusion pipeline that combines embedding similarity, token-level surprisal, and self-critique signals into continuous per-item exposure scores applicable to any benchmark without auxiliary data requirements.
    \item \textbf{Detection signal failure modes.} We document a directional reversal in Min-K\%++ at model scales below 1.5B parameters (Cohen's $d = -1.55$ to $-1.98$), where the signal is actively misleading rather than merely ineffective. This finding has immediate practical consequences for contamination detection at small scales.
    \item \textbf{Identifiability boundaries and exploratory domain analysis.} Simulation studies establish that reliable $\gamma$ recovery requires signal AUC $\geq 0.55$ and $J \geq 50$, a threshold well above current practice ($J = 8$--$30$). Exploratory per-domain $\gamma$ estimates at $J = 13$ reveal directional heterogeneity across MMLU cognitive domains (positive in STEM, negative in Social Sciences) but are largely underpowered, confirming the identifiability analysis and motivating future work with larger model pools.
\end{enumerate}